\documentclass[12pt,a4paper]{article}

\usepackage[utf8]{inputenc}
\usepackage[T1]{fontenc}
\usepackage[spanish,es-tabla]{babel}
\usepackage{amsmath,amssymb}
\usepackage{geometry}
\usepackage{pgfplots}
\pgfplotsset{compat=1.18}

\geometry{margin=2.5cm}

\title{Ejercicio a): Antitransformada de Laplace}
\author{}
\date{}

\begin{document}

\maketitle

\section*{Función dada}

Se considera:

\begin{equation}
    Y_1(s)=\frac{s-1}{(s+2)(s+4)}.
\end{equation}

\section*{Descomposición en fracciones simples}

Planteamos:

\begin{equation}
    \frac{s-1}{(s+2)(s+4)}
    =\frac{A}{s+2}+\frac{B}{s+4}.
\end{equation}

Multiplicando por $(s+2)(s+4)$:

\begin{align}
    s-1 &= A(s+4)+B(s+2),\\
    s-1 &= (A+B)s+(4A+2B).
\end{align}

Igualando coeficientes:

\begin{equation}
    \begin{cases}
        A+B=1,\\
        4A+2B=-1.
    \end{cases}
\end{equation}

Por lo tanto:

\begin{equation}
    A=-\frac{3}{2},\qquad B=\frac{5}{2}.
\end{equation}

Así,

\begin{equation}
    Y_1(s)=-\frac{3}{2}\frac{1}{s+2}
    +\frac{5}{2}\frac{1}{s+4}.
\end{equation}

\section*{Expresión en el dominio del tiempo}

Usando

\begin{equation}
    \mathcal{L}^{-1}\left\{\frac{1}{s+a}\right\}=e^{-at}u(t),
\end{equation}

se obtiene:

\begin{equation}
    \boxed{
    y_1(t)=\left(-\frac{3}{2}e^{-2t}+\frac{5}{2}e^{-4t}\right)u(t)
    }.
\end{equation}

Para $t\geq 0$:

\begin{equation}
    y_1(t)=-\frac{3}{2}e^{-2t}+\frac{5}{2}e^{-4t}.
\end{equation}

\section*{Deducción del cruce por cero}

Para encontrar el instante en que la señal se hace cero, igualamos la expresión a cero:

\begin{equation}
    \frac{5}{2}e^{-4t}-\frac{3}{2}e^{-2t}=0.
\end{equation}

Pasando el término negativo al segundo miembro:

\begin{equation}
    \frac{5}{2}e^{-4t}=\frac{3}{2}e^{-2t}.
\end{equation}

Multiplicando ambos lados por $2$:

\begin{equation}
    5e^{-4t}=3e^{-2t}.
\end{equation}

Dividiendo por $3e^{-4t}$, o equivalentemente multiplicando por $e^{4t}$:

\begin{align}
    \frac{5}{3}&=\frac{e^{-2t}}{e^{-4t}},\\
    \frac{5}{3}&=e^{-2t-(-4t)}=e^{2t}.
\end{align}

Aplicando el logaritmo neperiano a ambos lados:

\begin{align}
    \ln\left(\frac{5}{3}\right)&=\ln\left(e^{2t}\right),\\
    \ln\left(\frac{5}{3}\right)&=2t.
\end{align}

Finalmente, despejando $t$:

\begin{equation}
    \boxed{
    t=\frac{1}{2}\ln\left(\frac{5}{3}\right)
    \approx 0.2554
    }.
\end{equation}

\section*{Deducción del mínimo absoluto}

Definimos, para todo $t\in\mathbb{R}$,

\begin{equation}
    f(t)=\frac{5}{2}e^{-4t}-\frac{3}{2}e^{-2t}.
\end{equation}

Derivando con respecto a $t$:

\begin{equation}
    f'(t)=-10e^{-4t}+3e^{-2t}.
\end{equation}

Igualamos la derivada a cero:

\begin{align}
    -10e^{-4t}+3e^{-2t}&=0,\\
    3e^{-2t}&=10e^{-4t},\\
    e^{2t}&=\frac{10}{3}.
\end{align}

Aplicando el logaritmo neperiano:

\begin{equation}
    t_{\min}=\frac{1}{2}\ln\left(\frac{10}{3}\right)
    \approx 0.6020.
\end{equation}

La segunda derivada es

\begin{equation}
    f''(t)=40e^{-4t}-6e^{-2t}.
\end{equation}

En el punto crítico se cumple $e^{-2t_{\min}}=3/10$ y
$e^{-4t_{\min}}=9/100$. Por tanto,

\begin{equation}
    f''(t_{\min})=40\left(\frac{9}{100}\right)
    -6\left(\frac{3}{10}\right)=\frac{9}{5}>0,
\end{equation}

lo que confirma que el punto crítico es un mínimo. Evaluando la función:

\begin{align}
    f(t_{\min})&=\frac{5}{2}\left(\frac{9}{100}\right)
    -\frac{3}{2}\left(\frac{3}{10}\right)\\
    &=\boxed{-\frac{9}{40}=-0.225}.
\end{align}

Además, $f(t)\to+\infty$ cuando $t\to-\infty$ y
$f(t)\to0^-$ cuando $t\to+\infty$, por lo que este mínimo es absoluto.

\section*{Características y gráfica}

La señal comienza en $y_1(0^+)=1$. El cruce por cero ocurre en

\begin{equation}
    t=\frac{1}{2}\ln\left(\frac{5}{3}\right)\approx 0.2554,
\end{equation}

y tiende a cero desde valores negativos cuando $t\to\infty$.

\begin{center}
\begin{tikzpicture}
\begin{axis}[
    width=14cm,
    height=8cm,
    axis lines=middle,
    xlabel={$t$},
    ylabel={$y_1(t)$},
    xmin=0, xmax=3,
    ymin=-0.25, ymax=1.08,
    xtick={0,0.5,1,1.5,2,2.5,3},
    ytick={-0.2,0,0.2,0.4,0.6,0.8,1},
    grid=both,
    minor tick num=1,
    samples=200,
    domain=0:3,
    smooth,
    thick,
    blue
]
    \addplot {-(3/2)*exp(-2*x)+(5/2)*exp(-4*x)};
    \addplot[only marks, mark=*, mark size=1.6pt, red]
        coordinates {(0,1) (0.2554,0)};
\end{axis}
\end{tikzpicture}
\end{center}

\section*{Resultado final}

\begin{equation}
    \boxed{
    Y_1(s)=\frac{s-1}{(s+2)(s+4)}
    \quad\Longleftrightarrow\quad
    y_1(t)=\left(-\frac{3}{2}e^{-2t}+\frac{5}{2}e^{-4t}\right)u(t)
    }.
\end{equation}

\end{document}
